%& -shell-escape
% !TeX root = thesis-text.tex
% !TeX encoding = UTF-8
% !TeX program = XeLaTeX
%
% http://tex. stackexchange.com/questions/144\,908/how-can-i-make-a-reference-to-a-glossary-that-typesets-as-glossarytoctitle
%

\newglossaryentry{DoG}{
    name={DoG},
    description={
        difference of~gaussians; разница гауссианов;
        алгоритм выделения особенностей,
        выполняет вычитание одного размытой
        версии исходного изображения от~другого,
        менее размытой}
}

\newglossaryentry{PSIFT}{
    name={PSIFT},
    description={
        projective-scale-invariant feature transform;
        алгоритм поиска и~описания
        особых точек изображения. Устойчив по~отношению
        к~изменению масштаба, аффинным преобразованиям освещенности,
        изменению углов ориентации камеры
        и~проективный плоскости},
}

\newglossaryentry{SIFT}{
    name={SIFT},
    description={
        scale-invariant feature transform;
        алгоритм поиска и~описания
        особых точек изображения. Устойчив по~отношению
        к~изменению масштаба, аффинным преобразованиям
        освещенности \cite{Lowe:1999:SIFT}},
}

\newglossaryentry{vword}{
    name={Визуальное слово},
    description={
        Visual word; визуальное слово. Малый участок изображения или~кадра,
        содержащий какую~либо интересующую информацию
        в~каком~либо пространстве признаков
        (изменение цвета, изменения текстуры, и~т.~д.).
        Обычно представляет из~себя вектор пикселей}
}

\newglossaryentry{vkword}{
    name={Визуальное ключевое слово},
    description={
        Visual keyword; visual term; визуальное ключевое слово.
        Результат кластеризации
        участка изображения или~набора характеристик изображения
        в~пространстве признаков (центры кластеров)}
}

\newglossaryentry{embedding}{
    name={Вложение},
    description={
        специального вида отображение одного экземпляра
        некоторой математической структуры во~второй экземпляр такого~же~типа.
        Вложение некоторого объекта $X$ в $Y$ задаётся инъективным отображением,
        сохраняющим некоторую структуру},
}

\newglossaryentry{HammingEmbedding}{
    name=Вложение Хемминга,
    description={
        Hamming embedding; преобразование, при~котором точки
        некоторого $n$-мерного пространства (в~нашем случае видео),
        будут иметь сигнатуры близкие по~расстоянию Хемминга}
}

\newglossaryentry{HarrisDetektor}{
    name={Детектор Харриса},
    description={
        алгоритм выделения угловых точек на~изображении.
        Детектор Харриса инвариантен к~повороту 
        и~к~изменению интенсивности,
        но~не~инвариантен к~изменению масштаба}
}

\newglossaryentry{inverted-index}{
    name={Инвертированный индекс},
    description={
        inverted index; структура данных,
        в~которой для~каждого термина коллекции документов
        перечислены все документы в~коллекции,
        содержащие этот термин}
}

\newglossaryentry{frame}{
    name={Кадр},
    description={
        frame, фотографический кадр, кадрик;
        единичный элементарный фотоснимок на~киноплёнке,
        отдельная статическая картинка},
}

\newglossaryentry{keyframe}{
    name=Ключевой кадр,
    description={
    Существует два~различных значения термина.
    \begin{enumerate}
        \item   Опорный кадр; I-кадр. При~сжатии видео
                к~данным кадрам не~применяется межкадровое кодирование.
                Каждый опорный кадр сжимается независимо
                от~предыдущих и~следующих кадров.
        \item   Некоторый особый кадр видео,
                который выделяют для~каких~либо целей.
    \end{enumerate}
    В~литературе термин чаще всего используется в~первом значении.
    Мы~будем использовать термин во~втором смысле,
    если явно не~указано обратное}
}

\newglossaryentry{cosine-similarity}{
    name={Коэффициент Охаи},
    description={
        отношение общего числа признаков двух объектов
        к~среднему геометрическому двух объектов}
}

\newglossaryentry{LSH}{
    name={Локально чувствительное хеширование},
    description={
        locality-sensitive hashing; вероятностный метод понижения
        размерности многомерных данных.
        Идея состоит в~подборе хэш-функций,
        таких, что~для~похожих объектов значения этих~функций
        будут одинаковыми или~близкими по~некоторой метрике},
}

\newglossaryentry{bag-of-words}{
    name={Мешок слов},
    description={
        упрощенная модель данных,
        последовательность объектов представляется
        в~виде мешка (мультимножества),
        структура последовательности и~порядок объектов игнорируются}
}

\newglossaryentry{spoint}{
    name=Особая точка (кадра),
    description={
        ключевая точка; область изображения, в~которой градиент
        имеет локальные максимумы;
        наиболее заметные точки с~психофизической точки зрения;
        чаще всего в~качестве особых рассматривают угловые точки},
}

\newglossaryentry{recall}{
    name={Полнота (поиска)},
    description={
        $$
                \frac   {\abs{ \left\lbrace \text{\it найденные} \right\rbrace \cap  \left\lbrace \text{\it релевантные} \right\rbrace }}
                        {\abs{ \left\lbrace \text{\it релевантные} \right\rbrace }}
        $$ recall; отношение числа найденных релевантных объектов,
        к~общему числу релевантных объектов;
        неформально, можно сформулировать как
            <<как~много релевантных объектов было найдено
            по~отношению ко~всем релевантным>>}
}

\newglossaryentry{accuracy}{
    name={Правильность (поиска)},
    description={
        $$
                \frac   {\abs{ \left\lbrace \text{\it найденные} \right\rbrace }}
                        {\abs{ \left\lbrace \text{\it релевантные} \right\rbrace }}
        $$ accuracy; отношение количества найденных объектов
        к~количеству релевантных, т.~е. тех, которые должны были найти}
}

\newglossaryentry{HoughTransform}{
    name={Преобразование Хафа},
    description={
        численный алгоритм, используется для~поиска объектов,
        принадлежащих определённому классу фигур.
        Для~этого к~набору параметров применяется процедура голосования.
        Из~этого набора получаются объекты определённого класса}
}

\newglossaryentry{scene}{
    name=Сцена,
    description={scene;
    множество фотографических кадров,
    часть видео при~единстве места и~времени действия;
    состоит из~съемок и~выражает законченную идею}
}

\newglossaryentry{shot}{
    name=Съёмка,
    description={shot, кинематографический кадр, монтажный кадр, монтажный план;
    множество фотографических кадров,
    непрерывное действие между~пуском и~остановом камеры
    или~между~двумя монтажными склейками;
    часто называют <<сценой>>}
}

\newglossaryentry{precision}{
    name={Точность (поиска)},
    description={
        $$
                \frac   {\abs{ \left\lbrace \text{\it найденные} \right\rbrace \cap  \left\lbrace \text{\it релевантные} \right\rbrace }}
                        {\abs{ \left\lbrace \text{\it найденные} \right\rbrace }}
        $$ precision; отношение числа релевантных объектов, к~общему числу найденных объектов;
        неформально, можно сформулировать как
            <<как~много релевантных объектов
            из~того было найдено по~отношению ко~всем найденным>>}
}

\newglossaryentry{colormoment}{
    name={Цветовой момент},
    description={
        color moment; мера для~определения разницы между~изображениями
        по~их цветовым особенностям.
        Будем считать, что~цвета в~изображении распределены
        с~некоторой вероятностью.
        Моменты (среднее, дисперсия, третий момент) распределения
        могут быть использованы как~признаки для~идентификации
        этого~изображения}
}




